% GENERATED FILE. Edit canonical lessons, then run npm run generate:books.
% canonical-source-hash: fnv1a64:c4dd01b2a7e9a69f
% canonical-lessons: BN-C114-kuyasha, BN-C114-shishir, BN-C114-shikor, BN-C114-bij, BN-C114-kata

\chapter{Fog, Dew, Root, Seed, Thorn}
\label{ch:bn-fog-dew-root-seed-thorn}
\hlchaptermodality{114}

\begin{chapteropening}
\textbf{By the end of this chapter:} \emph{I can say fog, dew, a root, a seed and a thorn.}

Complete the last lesson of chapter 114: I can say fog, dew, a root, a seed and a thorn.
\end{chapteropening}

\section[\texorpdfstring{\bn{কুয়াশা} (kuyāshā)}{কুয়াশা (kuyāshā)}]{\bn{কুয়াশা} (kuyāshā) — fog}

\label{lesson:BN-C114-kuyasha}



\begin{quote}
Before the new one: say the Bengali for thunder, then the Bengali for a rainbow.
\end{quote}

\subsection*{You'll want to know: \bn{কুয়াশা}}
\textbf{\bn{কুয়াশা}} — \emph{kuyāshā} — \textquotedblleft{}fog\textquotedblright{}.

\begin{grammarlens}[title={What you've built}]
One more word for this chapter.
\end{grammarlens}

\subsection*{Your turn}
\begin{itemize}
\raggedright
  \item \emph{Say it:} \emph{kuyāshā}
  \item \emph{Say it:} \emph{kuyāshā}, once more
  \item \emph{Say it:} say \emph{kuyāshā}
  \item \emph{Recall:} say the Bengali for thunder, then the Bengali for a rainbow, then say \emph{kuyāshā} again
\end{itemize}

\begin{tcolorbox}[breakable,colback=teal!4,colframe=teal!35!black,title={Before you move on}]
What does \bn{কুয়াশা} mean? (\textquotedblleft{}Fog\textquotedblright{}.) Say it once more.
\end{tcolorbox}
\section[\texorpdfstring{\bn{শিশির} (shishir)}{শিশির (shishir)}]{\bn{শিশির} (shishir) — dew}

\label{lesson:BN-C114-shishir}



\begin{quote}
Before the new one: say the Bengali for a rainbow, then the Bengali for fog.
\end{quote}

\subsection*{You'll want to know: \bn{শিশির}}
\textbf{\bn{শিশির}} — \emph{shishir} — \textquotedblleft{}dew\textquotedblright{}.

\begin{grammarlens}[title={What you've built}]
One more word for this chapter.
\end{grammarlens}

\subsection*{Your turn}
\begin{itemize}
\raggedright
  \item \emph{Say it:} \emph{shishir}
  \item \emph{Say it:} \emph{shishir}, once more
  \item \emph{Say it:} say \emph{shishir}
  \item \emph{Recall:} say the Bengali for a rainbow, then the Bengali for fog, then say \emph{shishir} again
\end{itemize}

\begin{tcolorbox}[breakable,colback=teal!4,colframe=teal!35!black,title={Before you move on}]
What does \bn{শিশির} mean? (\textquotedblleft{}Dew\textquotedblright{}.) Say it once more.
\end{tcolorbox}
\section[\texorpdfstring{\bn{শিকড়} (shikôṛ)}{শিকড় (shikôṛ)}]{\bn{শিকড়} (shikôṛ) — a root}

\label{lesson:BN-C114-shikor}



\begin{quote}
Before the new one: say the Bengali for fog, then the Bengali for dew.
\end{quote}

\subsection*{You'll want to know: \bn{শিকড়}}
\textbf{\bn{শিকড়}} — \emph{shikôṛ} — \textquotedblleft{}a root\textquotedblright{}.

\begin{grammarlens}[title={What you've built}]
One more word for this chapter.
\end{grammarlens}

\subsection*{Your turn}
\begin{itemize}
\raggedright
  \item \emph{Say it:} \emph{shikôṛ}
  \item \emph{Say it:} \emph{shikôṛ}, once more
  \item \emph{Say it:} say \emph{shikôṛ}
  \item \emph{Recall:} say the Bengali for fog, then the Bengali for dew, then say \emph{shikôṛ} again
\end{itemize}

\begin{tcolorbox}[breakable,colback=teal!4,colframe=teal!35!black,title={Before you move on}]
What does \bn{শিকড়} mean? (\textquotedblleft{}A root\textquotedblright{}.) Say it once more.
\end{tcolorbox}
\section[\texorpdfstring{\bn{বীজ} (bij)}{বীজ (bij)}]{\bn{বীজ} (bij) — a seed}

\label{lesson:BN-C114-bij}



\begin{quote}
Before the new one: say the Bengali for dew, then the Bengali for a root.
\end{quote}

\subsection*{You'll want to know: \bn{বীজ}}
\textbf{\bn{বীজ}} — \emph{bij} — \textquotedblleft{}a seed\textquotedblright{}.

\begin{grammarlens}[title={What you've built}]
One more word for this chapter.
\end{grammarlens}

\subsection*{Your turn}
\begin{itemize}
\raggedright
  \item \emph{Say it:} \emph{bij}
  \item \emph{Say it:} \emph{bij}, once more
  \item \emph{Say it:} say \emph{bij}
  \item \emph{Recall:} say the Bengali for dew, then the Bengali for a root, then say \emph{bij} again
\end{itemize}

\begin{tcolorbox}[breakable,colback=teal!4,colframe=teal!35!black,title={Before you move on}]
What does \bn{বীজ} mean? (\textquotedblleft{}A seed\textquotedblright{}.) Say it once more.
\end{tcolorbox}
\section[\texorpdfstring{\bn{কাঁটা} (kā̃ṭā)}{কাঁটা (kā̃ṭā)}]{\bn{কাঁটা} (kā̃ṭā) — a thorn}

\label{lesson:BN-C114-kata}



\begin{quote}
Before the new one: say the Bengali for a root, then the Bengali for a seed.
\end{quote}

\subsection*{You'll want to know: \bn{কাঁটা}}
\textbf{\bn{কাঁটা}} — \emph{kā̃ṭā} — \textquotedblleft{}a thorn\textquotedblright{}.

\begin{grammarlens}[title={What you've built}]
One more word for this chapter.
\end{grammarlens}

\subsection*{Your turn}
\begin{itemize}
\raggedright
  \item \emph{Say it:} \emph{kā̃ṭā}
  \item \emph{Say it:} \emph{kā̃ṭā}, once more
  \item \emph{Say it:} say \emph{kā̃ṭā}
  \item \emph{Recall:} say the Bengali for a root, then the Bengali for a seed, then say \emph{kā̃ṭā} again
\end{itemize}

\begin{tcolorbox}[breakable,colback=teal!4,colframe=teal!35!black,title={Before you move on}]
What does \bn{কাঁটা} mean? (\textquotedblleft{}A thorn\textquotedblright{}.) Say it once more.
\end{tcolorbox}
